% ==== VII ====
\yoschapter{VII}{Dual Operators}
\label{chap:VII}


\yossection{1}{Dual Operators}
\label{sec:VII.1}

The notion of the transposed matrix may be extended to the notion of dual operator through 

\begin{theorem}\label{thm:VII.1.Theorem.1} Let $X$, $Y$ be locally convex linear topological spaces and $X_{s}^{\prime}, Y_{s}^{\prime}$ their strong dual spaces, respectively. Let $T$ be a linear operator on $D(T) \subseteq X$ into $Y$. Consider the points $\{x^{\prime}, y^{\prime}\}$ of the product space $X_{s}^{\prime} \times Y_{s}^{\prime}$ satisfying the condition 
\begin{equation}\label{eq:VII.1.1}
\langle T x, y ^{\prime} \rangle = \langle x, x ^{\prime} \rangle \quad \text{ for all } \quad x \in D (T).
\end{equation}
\end{theorem}

Then $x'$ is determined uniquely by $y'$ iff $D(T)$ is dense in $X$. 

\begin{proof}
 By the linearity of the problem, we have to consider the condition: 
\[
\langle x, x ^{\prime} \rangle = 0 \quad \text{ for all } \quad x \in D (T) \quad \text{ implies } \quad x ^{\prime} = 0. \]

Thus the ``if'' part is clear from the continuity of the linear functional $x'$. Assume that $D(T)^{a} \neq X$. Then there exists, by the Hahn-Banach theorem, an $x_{0}' \neq 0$ such that $\langle x, x_{0}' \rangle = 0$ for all $x \in D(T)$; consequently, the ``only if'' part is proved. 
\end{proof}

\begin{definition}\label{def:VII.1.Definition.1} $A$ linear operator $T'$ such that $T'y' = x'$ is defined through \eqref{eq:VII.1.1} iff $D(T)^a = X$. $T'$ is called the dual or conjugate operator of $T$; its domain $D(T')$ is the totality of those $y' \in Y_s'$ such that there exists $x' \in X_s'$ satisfying \eqref{eq:VII.1.1}, and $T'y' = x'$. Hence $T'$ is a linear operator defined on $D(T') \subseteq Y_s'$ into $X_s'$ such that 
\begin{equation}\label{eq:VII.1.2}
\langle T x, y ^{\prime} \rangle = \langle x, T ^{\prime} y ^{\prime} \rangle \quad \text{ for all } \quad x \in D (T) \quad \text{ and all } \quad y ^{\prime} \in D (T ^{\prime}).
\end{equation}
\end{definition}

\begin{theorem}\label{thm:VII.1.Theorem.2} If $D(T) = X$ and $T$ is continuous, then $T'$ is a continuous linear operator defined on $Y_s'$ into $X_s'$. 
\end{theorem}

\begin{proof}
 For any $y' \in Y_s'$, $\langle Tx, y' \rangle$ is a continuous linear functional of $x \in X$ and so there exists an $x' \in X_s'$ with $T'y' = x'$. Let $B$ be a bounded set of $X$. Then, by the continuity of $T$, the image $T \cdot B = \{Tx; x \in B\}$ is a bounded set of $Y$. Thus, by the defining relation $\langle Tx, y' \rangle = \langle x, x' \rangle$, the convergence to $0$ of $y'$ in the bounded convergence topology of $Y'$ (given in Chapter~\ref{chap:IV}, \ref{sec:IV.7}) implies the convergence of $x'$ in the bounded convergence topology of $X'$. Thus $T'$ is a continuous linear operator on $Y_s'$ into $X_s'$. 
\end{proof}

\begin{example}\label{ex:VII.1.Example.1} Let $X = Y$ be $n$-dimensional euclidean spaces normed by the $(l^2)$-norm. For any continuous linear operator $T \in L(X, X)$, set 
\[
T x = y, \quad \text{ where } \quad x = (x _{1}, x _{2}, \dots , x _{n}) \quad \text{ and } \quad y = (y _{1}, y _{2}, \dots , y _{n}).
\]
\end{example}

Then $y_{i} = \sum\limits_{j=1}^{n} t_{ij} x_{j} (i = 1, 2, \ldots, n)$ and so, for $z = (z_{1}, z_{2}, \ldots, z_{n})$, 
\[
\langle T x, z \rangle = \langle y, z \rangle = \sum\limits_{j} y _{j} z _{j} = \sum\limits_{i} \left(\sum\limits_{j} t _{ij} x _{j}\right) z _{i} = \sum\limits_{j} x _{j} \left(\sum\limits_{i} t _{ij} z _{i}\right).
\]

Thus $T'z = w$ is given by $w_{j} = \sum\limits_{i=1}^{n} t_{ij} z_{i} (j = 1, 2, \ldots, n)$. This proves that the matrix corresponding to $T'$ is the transposed matrix of the matrix corresponding to $T$. 

\begin{example}\label{ex:VII.1.Example.2} Let $X = Y$ be the real Hilbert space $(l^2)$, and let $T_n \in L(X, X)$ be defined by 
\[
T _{n} (x _{1}, x _{2}, \dots , x _{k}, \dots) = (x _{n}, x _{n+1}, x _{n+2}, \dots).
\]
\end{example}

Then from 
\[
\langle T _{n} (x _{1}, x _{2}, \dots), (z _{1}, z _{2}, \dots) \rangle = x _{n} z _{1} + x _{n+1} z _{2} + x _{n+2} z _{3} + \dots ,
\]
we obtain 
\[
T _{n} ^{\prime} (z _{1}, z _{2}, \dots) = (\overbrace {0 , 0 , \dots , 0}, z _{1}, z _{2}, \dots).
\]

Since $\|T_n(x_1, x_2, \ldots)\| = \left(\sum\limits_{m=n}^{\infty} x_m^2\right)^{1/2} \to 0$ as $n \to \infty$ and $\|T_n'(z_1, z_2, \ldots)\|$ $= \|(z_1, z_2, \ldots)\|$, we have 

\begin{proposition}\label{prop:VII.1.Proposition.1} The mapping $T \to T'$ of $L(X, Y)$ into $L(Y_s', X_s')$ is not, in general, continuous in the simple convergence topology of operators, that is, $\lim\limits_{n \to \infty} T_n x = Tx$ for all $x \in X$ does not necessarily imply $\lim\limits_{n \to \infty} T_n' y' = T'y'$ for all $y' \in Y'$, in the strong dual topology of $X_s'$. 
\end{proposition}

\begin{theoremn}[Theorem~\ref{thm:VII.1.Theorem.2}']\yosnamedtag{2'}\label{thm:VII.1.Theorem.2prime} Let $T$ be a bounded linear operator on a normed linear space $X$ into a normed linear space $Y$. Then the dual operator $T'$ is a bounded linear operator on $Y_s'$ into $X_s'$ such that 
\begin{equation}\label{eq:VII.1.3}
\| T \| = \left\| T ^{\prime} \right\|.
\end{equation}
\end{theoremn}

\begin{proof}
 From the defining relation $\langle Tx, y'\rangle = \langle x, x'\rangle$, we obtain 
\[
\begin{aligned}
&\left\| T ^{\prime} y ^{\prime} \right\| = \left\| x ^{\prime} \right\| = \sup\limits _{\| x \| \leqq 1} \left| \langle x, x ^{\prime} \rangle \right| = \sup\limits _{\| x \| \leqq 1} \left| \langle T x, y ^{\prime} \rangle \right| \leqq \left\| y ^{\prime} \right\| \cdot \sup\limits _{\| x \| \leqq 1} \left\| T x \right\| \leqq \left\| y ^{\prime} \right\| \cdot \left\| T \right\|,
\end{aligned}
\]
and so $\| T' \| \leqq \| T \|$. The reverse inequality is proved as follows. For any $x_0 \in X$, there exists an $f_0 \in Y'$ such that $\| f_0 \| = 1$ and $f_0(Tx_0) = \langle Tx_0, f_0 \rangle = \| Tx_0 \|$. Thus $f_0' = T'f_0$ satisfies $\langle x_0, f_0' \rangle = \| Tx_0 \|$ and so 
\begin{equation}\left\| T x _{0} \right\| = \left\langle x _{0}, T ^{\prime} f _{0} \right\rangle \leqq \left\| T ^{\prime} \right\| \left\| f _{0} \right\| \cdot \left\| x _{0} \right\| = \left\| T ^{\prime} \right\| \cdot \left\| x _{0} \right\|, i. e., \left\| T \right\| \leqq \left\| T ^{\prime} \right\|. \notag\end{equation}
\end{proof}

\begin{theorem}\label{thm:VII.1.Theorem.3} \begin{enumerate}[i)]
\item If $T$ and $S$ are $\in L(X, Y)$, then $(\alpha T + \beta S)' = (\alpha T' + \beta S')$.
\item Let linear operators $T, S$ be such that $D(T), D(S), R(T)$ and $R(S)$ are all contained in $X$. If $S$ is $\in L(X, X)$ and $D(T)^a = X$, then 
\begin{equation}\label{eq:VII.1.4}
(S T) ^{\prime} = T ^{\prime} S ^{\prime}.
\end{equation}
\end{enumerate}
\end{theorem}

If, moreover, $D(TS)^{a} = X$, then 
\begin{equation}\label{eq:VII.1.5}
(T S) ^{\prime} \supseteq S ^{\prime} T ^{\prime}, \text{ i.e., } (T S) ^{\prime} \text{ is an extension of } S ^{\prime} T ^{\prime}.
\end{equation}

\begin{proof}
 \begin{enumerate}[i)]
\item is clear.
\item $D(ST) = D(T)$ is dense in $X$, and so $(ST)'$ exists. If $y \in D((ST)')$, then, for any $x \in D(T) = D(ST)$, $\langle Tx, S'y \rangle = \langle STx, y \rangle = \langle x, (ST)'y \rangle$. This shows that $S'y \in D(T')$ and $T'S'y = (ST)'y$, that is, $(ST)' \subseteq T'S'$. Let, conversely, $y \in D(T'S')$, i.e., $S'y \in D(T')$. Then, for any $x \in D(T) = D(ST)$, $\langle STx, y \rangle = \langle Tx, S'y \rangle = \langle x, T'S'y \rangle$. This shows that $y \in D((ST)')$ and $(ST)'y = T'S'y$, that is, $T'S' \subseteq (ST)'$. We have thus proved \eqref{eq:VII.1.4}.
\end{enumerate}



To prove \eqref{eq:VII.1.5}, let $y \in D(S'T') = D(T')$. Then, for any $x \in D(TS)$, $\langle TSx, y \rangle = \langle Sx, T'y \rangle = \langle x, S'T'y \rangle$. This shows that $y \in D((TS)')$ and $(TS)'y = S'T'y$, that is, $S'T' \subseteq (TS)'$.
\end{proof}

\yossection{2}{Adjoint Operators}
\label{sec:VII.2}

The notion of transposed conjugate matrix may be extended to the notion of adjoint operator through 

\begin{definition}\label{def:VII.2.Definition.1} Let $X$, $Y$ be Hilbert spaces, and $T$ a linear operator defined on $D(T) \subseteq X$ into $Y$. Let $D(T)^{a} = X$ and let $T'$ be the dual operator of $T$. Thus $\langle Tx, y' \rangle = \langle x, T'y' \rangle$ for $x \in D(T)$, $y' \in D(T')$. If we denote by $J_{X}$ the one-to-one norm-preserving conjugate linear correspondence $X_{s}' \ni f \leftrightarrow y_{f} \in X$ (defined in Corollary~\ref{cor:III.6.Corollary.1} in Chapter~\ref{chap:III}, \ref{sec:III.6}), then $\langle Tx, y' \rangle = y'(Tx) = (Tx, J_{Y}y')$, $\langle x, T'y' \rangle = (T'y')(x) = (x, J_{X}T'y')$. 
\end{definition}

We have thus 
\[
(T x, J _{Y} y ^{\prime}) = (x, J _{X} T ^{\prime} y ^{\prime}), \text{ that is, } (T x, y) = (x, J _{X} T ^{\prime} J _{Y} ^{-1} y).
\]

In the special case when $Y = X$, we write 
\[
T ^{*} = J _{X} T ^{\prime} J _{X} ^{-1}
\]
and call $T^*$ the adjoint operator of $T$. 

\begin{remarku} If $X$ is the complex Hilbert space $(l^{2})$, we see, as in the Example in the preceding section, that the matrix corresponding to $T^{*}$ is the transposed conjugate matrix of the matrix corresponding to $T$. 
\end{remarku}

As in the case of dual operators, we can prove 

\begin{theorem}\label{thm:VII.2.Theorem.1} $T^*$ exists iff $D(T)^a = X$. It is defined as follows: $y \in X$ is in the domain of $D(T^*)$ iff there exists a $y^* \in X$ such that 
\begin{equation}\label{eq:VII.2.1}
(T x, y) = (x, y ^{*}) \quad \text{ holds for all } \quad x \in D (T),
\end{equation}
and we define $T^{*}y = y^{*}$. 
\end{theorem}

We can rewrite the above theorem in terms of the graph $G(A)$ of the linear operator $A$ (the graph was introduced in Chapter~\ref{chap:II}, \ref{sec:II.6}): 

\begin{theorem}\label{thm:VII.2.Theorem.2} We introduce a continuous linear operator $V$ on $X \times X$ into $X \times X$ by 
\begin{equation}\label{eq:VII.2.2}
V \{x, y \} = \{- y, x \}.
\end{equation}
\end{theorem}

Then $(VG(T))^{\perp}$ is the graph of a linear operator iff $D(T)^{a} = X$, and, in fact, we have 
\begin{equation}\label{eq:VII.2.3}
G (T ^{*}) = (V G (T)) ^{\perp}.
\end{equation}

\begin{proof}
 The condition $\{-Tx, x\} \perp \{y, y^*\}$ is equivalent to $(Tx, y) = (x, y^*)$. Thus Theorem~\ref{thm:VII.2.Theorem.2} is proved by Theorem~\ref{thm:VII.2.Theorem.1}. 
\end{proof}

\begin{corollaryu} $T^{*}$ is a closed linear operator, since the orthogonal complement of a linear subspace is a closed linear subspace. 
\end{corollaryu}

\begin{theorem}\label{thm:VII.2.Theorem.3} Let $T$ be a linear operator on $D(T) \subseteq X$ into $X$ such that $D(T)^a = X$. Then $T$ admits a closed linear extension iff $T^{**} = (T^*)^*$ exists, i.e., iff $D(T^*)^a = X$. 
\end{theorem}

\begin{proof}
 Sufficiency. We have $T^{**} \supseteq T$ by definition, and $T^{**} = (T^{*})^{*}$ is closed by the above Corollary. 

Necessity. Let $S$ be a closed extension of $T$. Then $G(S)$ contains $G(T)^{a}$ as a closed linear subspace, and so $G(T)^{a}$ is the graph of a linear operator. But $G(T)^{a} = G(T)^{\perp\perp} = (G(T)^{\perp})^{\perp}$ by the continuity of the scalar product, and, moreover, by $VG(T^{*}) = G(T)^{\perp}$, we obtain $(VG(T^{*}))^{\perp} = G(T)^{\perp\perp}$. Therefore, $(VG(T^{*}))^{\perp}$ is the graph of a linear operator. Thus by Theorem~\ref{thm:VII.2.Theorem.2} $D(T^{*})^{a} = X$ and $T^{**}$ exists. 
\end{proof}

\begin{corollaryu} Under the condition that $D(T)^{a}=X$, $T$ is closed linear iff $T=T^{**}$. 
\end{corollaryu}

\begin{proof}
 The sufficiency is clear. Necessity is proved by observing the formula $G(T)^{a} = G(T^{**})$ obtained above. For, $G(T) = G(T)^{a}$ implies that $T = T^{**}$. 
\end{proof}

\begin{theorem}\label{thm:VII.2.Theorem.4} An everywhere defined closed linear operator $T$ is a continuous linear operator. 
\end{theorem}

\begin{proof}
 Clear from the closed graph theorem. 
\end{proof}

\begin{theorem}\label{thm:VII.2.Theorem.5} If $T$ is a bounded linear operator, then $T^*$ is also bounded linear and 
\begin{equation}\label{eq:VII.2.4}
\left\| T \right\| = \left\| T ^{*} \right\|.
\end{equation}
\end{theorem}

\begin{proof}
 Similar to the case of dual operators. 
\end{proof}


\yossection{3}{Symmetric Operators and Self-adjoint Operators}
\label{sec:VII.3}

A Hermitian matrix is a matrix which is equal to its transposed conjugate matrix. It is known that such a matrix can be transformed into a diagonal matrix by a suitable (complex) rotation of the vector space on which the matrix operates as a linear operator. The notion of the Hermitian matrices is extended to the notion of self-adjoint operators in a Hilbert space. 

\begin{definition}\label{def:VII.3.Definition.1} Let $X$ be a Hilbert space. A linear operator on $D(T) \subseteq X$ into $X$ is called symmetric if $T^* \supseteq T$, i.e., if $T^*$ is an extension of $T$. Note that the condition of the existence of $T^*$ implies that $D(T)^a = X$. 
\end{definition}

\begin{proposition}\label{prop:VII.3.Proposition.1} If $T$ is symmetric, then $T^{**}$ is also symmetric. 
\end{proposition}

\begin{proof}
 Since $T$ is symmetric, we have $D(T^{*}) \supseteq D(T)$ and $D(T)^{a} = X$. Hence $D(T^{*})^{a} = X$ and so $T^{**} = (T^{*})^{*}$ exists. $T^{**}$ is surely an extension of $T$ and so $D(T^{**}) \supseteq D(T)$. Thus, again by $D(T)^{a} = X$, we have $D(T^{**})^{a} = X$ and so $T^{***} = (T^{**})^{*}$ exists. We have, from $T^{*} \supseteq T$, $T^{**} \subseteq T^{*}$ and hence $T^{***} \supseteq T^{**}$ which proves that $T^{**}$ is symmetric. Corollary. A symmetric operator $T$ has a closed symmetric extension $T^{**} = (T^{*})^{*}$. 
\end{proof}

\begin{definition}\label{def:VII.3.Definition.2} $A$ linear operator $T$ is called self-adjoint if $T = T^{*}$. 
\end{definition}

\begin{proposition}\label{prop:VII.3.Proposition.2} $A$ self-adjoint operator is closed. An everywhere defined symmetric operator is bounded and self-adjoint. 
\end{proposition}

\begin{proof}
 Being the adjoint of itself, a self-adjoint operator is closed. The last assertion is proved by the fact that an everywhere defined closed operator is bounded (closed graph theorem). 
\end{proof}

\begin{example}[integral operator of the Hilbert-Schmidt type]\label{ex:VII.3.Example.1} Let $-\infty \leqq a < b \leqq \infty$ and consider $L^2(a, b)$. Let $K(s, t)$ be a complexvalued measurable function for $a \leqq s, t \leqq b$ such that 
\[
\int_{a} ^{b} \int_{a} ^{b} | K (s, t) | ^{2} \dd s \dd t < \infty .
\]
\end{example}

For any $x(t) \in L^{2}(a, b)$, we define the operator $K$ by 
\[
(K \cdot x) (s) = \int_{a} ^{b} K (s, t) x (t) \dd t.
\]

We have, by Schwarz' inequality and the Fubini-Tonelli theorem. 
\[
\int_{a} ^{b} | (K \cdot x) (s) | ^{2} \dd s \leqq \int_{a} ^{b} \int_{a} ^{b} | K (s, t) | ^{2} \dd t \dd s \int_{a} ^{b} | x (t) | ^{2} \dd t.
\]

Hence $K$ is a bounded linear operator on $L^2(a, b)$ into $L^2(a, b)$ such that $\|K\| \leqq \left( \int_{a}^{b} \int_{a}^{b} |K(s, t)|^2 \dd s \dd t \right)^{1/2}$. It is easy to see that the operator $K^*$ is defined by $(K^* y)(t) = \int_{a}^{b} \overline{K(t, s)} y(s) \dd s$. Hence $K$ is self-adjoint iff $K(s, t) = \overline{K(t, s)}$ for a.e. $s, t$. 

\begin{example}[the coordinate operator in quantum mechanics]\label{ex:VII.3.Example.2} Let $X = L^2(-\infty, \infty)$. Let $D = \{x(t); x(t) \text{ and } t \cdot x(t) \text{ both } \in L^2(-\infty, \infty)\}$. Then the operator $T$ defined by $Tx(t) = t \cdot x(t)$ on $D$ is self-adjoint. 
\end{example}

\begin{proof}
 It is clear that $D^a = X$, since the linear combinations of defining functions of finite intervals are strongly dense in $L^2(-\infty, \infty)$. Let $y \in D(T^*)$ and set $T^*y = y^*$. Then, for all $x \in D = D(T)$, 
\[
\int_{- \infty} ^{\infty} t x (t) \overline {{y (t)}} \dd t = \int_{- \infty} ^{\infty} x (t) \overline {{y ^{*} (t)}} \dd t. \]

If we take for $x(t)$ the defining function of the interval $[\alpha, t_0]$, we have $\int_{\alpha}^{t_0} t \cdot \overline{y(t)} \dd t = \int_{\alpha}^{t_0} y^*(t) \dd t$, and hence, by differentiation, $t_0 \cdot \overline{y(t_0)} = \overline{y^*(t_0)}$ for a.e. $t_0$. Thus $y \in D$ and $T^* y(t) = t \cdot y(t)$. Conversely, it is clear that $y \in D$ implies that $y \in D(T^*)$ and $T^* y(t) = t \cdot y(t)$. 
\end{proof}

\begin{example}[the momentum operator in quantum mechanics]\label{ex:VII.3.Example.3} Let $X = L^{2}(-\infty, \infty)$. Let $D$ be the totality of $x(t) \in L^{2}(-\infty, \infty)$ such that $x(t)$ is absolutely continuous on every finite interval with the derivative $x'(t) \in L^{2}(-\infty, \infty)$. Then the operator $T$ defined by $Tx(t) = i^{-1}x'(t)$ on $D$ is self-adjoint. 
\end{example}

\begin{proof}
 Let a continuous function $x_{n}(t)$ be defined by 
\begin{equation}x _{n} (t) = 1 \text{ for } t \in [ \alpha , t _{0} ], x _{n} (t) = 0 \text{ for } t \leqq \alpha - n ^{-1} \text{ and for } t \geqq t _{0} + n ^{-1}, \notag\end{equation}
$x_{n}(t)$ is a linear function on $[\alpha - n^{-1}, \alpha]$ and on $[t_0, t_0 + n^{-1}]$. 

Then the linear combinations of functions of the form $x_{n}(t)$ with different values of $\alpha, t_0$ and $n$ are dense in $L^2(-\infty, \infty)$. Thus $D$ is dense in $X$. Let $y \in D(T^*)$ and $T^*y = y^*$. Then for any $x \in D$, 
\[
\int_{- \infty} ^{\infty} i ^{-1} x ^{\prime} (t) \overline {{{y (t)}}} \dd t = \int_{- \infty} ^{\infty} x (t) \overline {{{y ^{*} (t)}}} \dd t.
\]

If we take $x_{n}(t)$ for $x(t)$, we obtain 
\[
n \int_ {x - n ^{-1}} ^{\alpha} i ^{-1} \overline {{{y (t)}}} \dd t - n \int_ {t _{0}} ^ {t _{0} + n ^{-1}} i ^{-1} \overline {{{y (t)}}} \dd t = \int_{- \infty} ^{\infty} x _{n} (t) y ^{*} (t) \dd t,
\]
and so, by letting $n \to \infty$, we obtain $i^{-1}(\overline{y(\alpha)} - \overline{y(t_0)}) = \int_{\alpha}^{t_0} \overline{y^*(t)} \dd t$ for a.e. $\alpha$ and $t_0$. It is clear, by Schwarz' inequality, that $y^*(t)$ is integrable over any finite interval. Thus $y(t_0)$ is absolutely continuous in $t_0$ over any finite interval, and so we have $i^{-1}y'(t_0) = y^*(t_0)$ for a.e. $t_0$. Hence $y \in D$ and $T^*y(t) = i^{-1}y'(t)$. Let, conversely, $y \in D$. Then, by partial integration, 
\[
\int_{a} ^{b} i ^{-1} x ^{\prime} (t) \overline {{{y (t)}}} \dd t = i ^{-1} [ x (t) \overline {{{y (t)}}} ] _{a} ^{b} + \int_{a} ^{b} x (t) \overline {{{(i ^{-1} y ^{\prime} (t))}}} \dd t.
\]

By the integrability of $x(t) \overline{y(t)}$ over $(-\infty, \infty)$, we see that $\lim\limits_{a \downarrow -\infty, b \uparrow \infty} |[x(t)y(t)]_a^b| = 0$, and so $\int_{-\infty}^{\infty} i^{-1} x'(t) \overline{y(t)} \dd t = \int_{-\infty}^{\infty} x(t) (\overline{i^{-1} y'(t)}) \dd t$. Thus $y \in D(T^*)$ and $T^* y(t) = i^{-1} y'(t)$. 
\end{proof}

\begin{theorem}\label{thm:VII.3.Theorem.1} If a self-adjoint operator $T$ admits the inverse $T^{-1}$, then $T^{-1}$ is also self-adjoint. 
\end{theorem}

\begin{proof}
 $T = T^{*}$ is equivalent to $(VG(T))^{\perp} = G(T)$. We have also $G(T^{-1}) = VG(-T)$. Hence, by $(-T)^{*} = -T^{*} = -T$, $(VG(-T))^{\perp} = G(-T)$ and so 
\[
(V G (T ^{-1})) ^{\perp} = G (- T) ^{\perp} = (V G (- T)) ^{\perp \perp} = V G (- T) = G (T ^{-1}),
\]
that is, $(T^{-1})^{*} = T^{-1}$. 


We have used, in the above proof, the fact that $(VG(-T))^a = VG(-T)$ in virtue of the closedness of $(-T)$. 
\end{proof}

\begin{corollaryu} $A$ symmetric operator $T$ in a Hilbert space $X$ is self-adjoint if $D(T) = X$ or if $R(T) = X$. 
\end{corollaryu}

\begin{proof}
 The case $D(T) = X$ was proved already. We shall prove the case $R(T) = X$. $Tx = 0$ implies $0 = (Tx, y) = (x, Ty)$ for all $y \in D(T)$, and so, by $R(T) = X$, we must have $x = 0$. Therefore the inverse $T^{-1}$ exists which is surely symmetric with $T$. $D(T^{-1}) = R(T) = X$, and so the everywhere defined symmetric operator $T^{-1}$ must be self-adjoint. Hence $T = (T^{-1})^{-1}$ is self-adjoint by Theorem~\ref{thm:VII.3.Theorem.1}. 
\end{proof}

We can construct self-adjoint operators from a closed linear operator. More precisely, we have 

\begin{theorem}[J. von \textsc{Neumann} {\cite{neumann232}}]\label{thm:VII.3.Theorem.2} For any closed linear operator $T$ in a Hilbert space $X$ such that $D(T)^{a} = X$, the operators $T^{*}T$ and $IT^{*}$ are self-adjoint, and $(I + T^{*}T)$ and $(I + TT^{*})$ both admit bounded linear inverses. 
\end{theorem}

\begin{proof}
 We know that, in the product space $X \times X$, $G(T)$ and $VG(T^{*})$ are closed linear subspaces orthogonal to each other and spanning the whole product space $X \times X$. Hence, for any $h \in X$, we have the uniquely determined decomposition 
\begin{equation}\label{eq:VII.3.1}
\{h, 0 \} = \{x, T x \} + \{- T ^{*} y, y \} \quad \text{ with } \quad x \in D (T), y \in D (T ^{*}). \end{equation}

Thus $h = x - T^{*}y$, $0 = Tx + y$. Therefore 
\begin{equation}\label{eq:VII.3.2}
x \in D (T ^{*} T) \quad \text{ and } \quad x + T ^{*} T x = h.
\end{equation}

Because of the uniqueness of decomposition \eqref{eq:VII.3.1}, $x$ is uniquely determined by $h$, and so the everywhere defined inverse $(I + T^{*}T)^{-1}$ exists. 

For any $h, k \in X$, let 
\[
x = (I + T ^{*} T) ^{-1} h, y = (I + T ^{*} T) ^{-1} k.
\]

Then $x$ and $y \in D(T^*T)$ and, by the closedness of $T$, $(T^*)^* = T$. Hence 
\[
\begin{aligned}
&(h, (I + T ^{*} T) ^{-1} k)  = ((I + T ^{*} T) x, y) = (x, y) + (T ^{*} T x, y) \\
&= (x, y) + (T x, T y) = (x, y) + (x, T ^{*} T y) = (x, (I + T ^{*} T) y) = ((I + T ^{*} T) ^{-1} h, k),
\end{aligned}
\]
which proves that the operator $(I + T^{*}T)^{-1}$ is self-adjoint. As an everywhere defined self-adjoint operator, $(I + T^{*}T)^{-1}$ is a bounded operator. By Theorem~\ref{thm:VII.3.Theorem.1}, its inverse $(I + T^{*}T)$ and hence $T^{*}T$ are self-adjoint. 

Since $T$ is closed, we have $(T^{*})^{*} = T$, and so, by what was proved above, $TT^{*} = (T^{*})^{*}T^{*}$ is self-adjoint and $(I + TT^{*})$ has a bounded linear inverse. 
\end{proof}

We next give an example of a non-self-adjoint, symmetric operator: 

\begin{example}\label{ex:VII.3.Example.4} Let $X = L^2(0,1)$. Let $D$ be the totality of absolutely continuous functions $x(t) \in L^2(0,1)$ such that $x(0) = x(1) = 0$ and $x'(t) \in L^2(0,1)$. Then the operator $T_1$ defined by $T_1x(t) = i^{-1}x'(t)$ on $D = D(T_1)$ is symmetric but not self-adjoint. 
\end{example}

\begin{proof}
 We shall prove that $T_1^* = T_2$, where $T_2$ is defined by: 
\[
T _{2} x (t) = i ^{-1} x ^{\prime} (t) \quad \text{ on } \quad D (T _{2}) = \{x (t) \in L ^{2} (0, 1); x (t) \text{ is }
\]
absolutely continuous such that $x'(t) \in L^2(0, 1)$. 

Since $D = D(T_{1})$ is dense in $L^2(0,1)$, the operator $T_1^*$ is defined. Let $y \in D(T_1^*)$ and set $T_1^* y = y^*$. Then, for any $x \in D = D(T_1)$, 
\[
\int_{0} ^{1} i ^{-1} x ^{\prime} (t) \overline {{y (t)}} \dd t = \int_{0} ^{1} x (t) \overline {{y ^{*} (t)}} \dd t.
\]

By partial integration, we obtain, remembering $x(0) = x(1) = 0$, 
\[
\int_{0} ^{1} x (t) \overline {{{y ^{*} (t)}}} \dd t = - \int_{0} ^{1} x ^{\prime} (t) \overline {{{Y ^{*} (t)}}} \dd t, \text{ where } Y ^{*} (t) = \int_{0} ^{t} y ^{*} (s) \dd s.
\]

Hence, by $x(1) = \int_{0}^{1} x'(t) \dd t = 0$, we have, for any constant $c$, 
\[
\int_{0} ^{1} x ^{\prime} (t) \overline {{{(Y ^{*} (t))}}} - i ^{-1} \overline {{{y (t)}}} - \overline {{{c}}}) \dd t = 0 \quad \text{ for all } \quad x \in D (T _{1}).
\]

On the other hand, for any $z(t) \in L^2(0, 1)$, the function $Z(t) = \int_0^t z(t) \dd t - t \int_0^1 z(t) \dd t$ surely belongs to $D(T_1)$. Hence, taking $Z(t)$ for the above $x(t)$, we obtain 
\[
\int_{0} ^{1} \left\{z (t) - \int_{0} ^{1} z (t) \dd t \right\} \cdot \overline {{(Y ^{*} (t)}} - i ^{-1} y (t) - \bar {c}) \dd t = 0.
\]

If we take the constant $c$ in such a way that $\int_0^1 (Y^* (t) - i^{-1}y(t) - c)\dd t = 0$, then 
\[
\int_{0} ^{1} z (t) \overline {{(Y ^{*} (t)}} - \overline {{i ^{-1} y (t)}} - \overline {{c}}) \dd t = 0,
\]
and so, by the arbitrariness of $z \in L^2(0,1)$, we must have $Y^*(t) = \int_0^t y^*(t) \dd t = i^{-1} y(t) + c$. Hence $y \in D(T_2)$ and $T_2y = y^*$. This proves that $T_1^* \subseteq T_2$. It is also clear, by partial integration, that $T_2 \subseteq T_1^*$ and so $T_2 = T_1^*$. 
\end{proof}

\begin{theorem}\label{thm:VII.3.Theorem.3} If $H$ is a bounded self-adjoint operator, then 
\begin{equation}\label{eq:VII.3.3}
\left\| H \right\| = \sup\limits _{\left\| x \right\| \leqq 1} \left| (H x, x) \right|.
\end{equation}
\end{theorem}

\begin{proof}
 Set $\sup\limits_{\|x\|\leqq 1}|(Hx,x)| = \gamma$. Then, by $|(Hx,x)|\leqq \|Hx\|\|x\|,$ $\gamma \leqq \|H\|$. For any real number $\lambda$, we have 
\begin{equation}\left| (H (y \pm \lambda z), y \pm \lambda z) \right| = \left| (H y, y) \pm 2 \lambda \operatorname{Re} (H y, z) + \lambda^{2} (H z, z) \right| \leqq \gamma \| y \pm \lambda z \| ^{2}. \notag\end{equation}

Hence 
\[
\left| 4 \lambda \operatorname{Re} (H y, z) \right| \leqq \gamma (\left\| y + \lambda z \right\| ^{2} + \left\| y - \lambda z \right\| ^{2}) = 2 \gamma (\left\| y \right\| ^{2} + \lambda^{2} \left\| z \right\| ^{2}).
\]

By taking $\lambda = \|y\|/\|z\|$, we obtain $|\operatorname{Re}(Hy,z)| \leqq \gamma \|y\|/\|z\|$. Hence, by substituting $ze^{i\theta}$ for $z$, we obtain $|(Hy,z)| \leqq \gamma \|y\|/\|z\|$, and so 
\[
(H y, H y) = \| H y \| ^{2} \leqq \gamma \| y \| \| H y \|, \text{ i.e., } \| H \| \leqq \gamma .
\]
\end{proof}


\yossection{4}{Unitary Operators. The Cayley Transform}
\label{sec:VII.4}

A symmetric operator is not necessarily a bounded operator. Various investigations of a symmetric operator $H$ may be made through the continuous operator $(H - iI)(H + iI)^{-1}$ called the Cayley transform of $H$. We shall begin with the notion of isometric operators. 

\begin{definition}\label{def:VII.4.Definition.1} $A$ bounded linear operator $T$ on a Hilbert space $X$ into $X$ is called (bounded) isometric if $T$ leaves the scalar product invariant: 
\begin{equation}\label{eq:VII.4.1}
(T x, T y) = (x, y) \quad \text{ for all } \quad x, y \in X.
\end{equation}
\end{definition}

If, in particular, $R(T) = X$, then a (bounded) isometric operator $T$ is called a unitary operator. 

\begin{proposition}\label{prop:VII.4.Proposition.1} For a bounded linear operator $T$, condition \eqref{eq:VII.4.1} is equivalent to the condition of the isometry 
\begin{equation}\label{eq:VII.4.2}
\left\| T x \right\| = \left\| x \right\| \quad \text{ for all } \quad x \in X.
\end{equation}
\end{proposition}

\begin{proof}
 It is clear that \eqref{eq:VII.4.1} implies \eqref{eq:VII.4.2}. We have, by \eqref{eq:VII.4.2}, 
\[ \begin{aligned}
&4 \operatorname{Re} (T x, T y)  = \| T (x + y) \| ^{2} - \| T (x - y) \| ^{2} = \| x + y \| ^{2} - \| x - y \| ^{2} = 4 \operatorname{Re} (x, y).
\end{aligned}
\]

By taking $iy$ in place of $y$, we also obtain $4\operatorname{Im}(Tx, Ty) = 4\operatorname{Im}(x, y)$, and so \eqref{eq:VII.4.2} implies \eqref{eq:VII.4.1}. 
\end{proof}

\begin{proposition}\label{prop:VII.4.Proposition.2} $A$ bounded linear operator on a Hilbert space $X$ into $X$ is unitary iff $T^{*} = T^{-1}$. 
\end{proposition}

\begin{proof}
 If $T$ is unitary, then $T^{-1}$ surely exists in virtue of condition \eqref{eq:VII.4.2}, and $D(T^{-1}) = R(T) = X$. Moreover, by \eqref{eq:VII.4.1}, $T^{*}T = I$ and so $T^{*} = T^{-1}$. Conversely, the condition $T^{*} = T^{-1}$ implies $T^{*}T = I$ which is the condition of the invariance of the scalar product. Moreover, $T^{*} = T^{-1}$ implies that $R(T) = D(T^{-1}) = D(T^{*}) = X$ and hence we see that $T$ must be unitary. 
\end{proof}

\begin{example}\label{ex:VII.4.Example.1} Let $X = L^{2}(-\infty, \infty)$. Then, for any real number a, the operator $T$ defined by $Tx(t) = x(t + a)$ on $L^{2}(-\infty, \infty)$ is unitary. 
\end{example}

\begin{example}\label{ex:VII.4.Example.2} The Fourier transform on $L^2(\mathbb{R}^{n})$ onto $L^2(\mathbb{R}^{n})$ is unitary, since it leaves the scalar product $(f, g) = \int f(x) \overline{g(x)} \dd x$ invariant. 
\end{example}

\begin{definition}\label{def:VII.4.Definition.2} Let $X$ be a Hilbert space. A linear operator $T$ defined on $D(T) \subseteq X$ into $X$ such that $D(T)^a = X$ is called normal if 
\begin{equation}\label{eq:VII.4.3}
T T ^{*} = T ^{*} T.
\end{equation}
\end{definition}

Self-adjoint operators and unitary operators are normal. 

\yossub{The Cayley Transform}

\begin{theorem}[J. \textsc{von} \textsc{Neumann} {\cite{neumann228}}]\label{thm:VII.4.Theorem.1} Let $H$ be a closed symmetric operator in a Hilbert space $X$. Then the continuous (but not necessarily everywhere defined) inverse $(H + iI)^{-1}$ exists, and the operator 
\begin{equation}\label{eq:VII.4.4}
U _{H} = (H - iI) (H + iI) ^{-1} \text{ with the domain } D(U_{H}) = D((H + iI) ^{-1})
\end{equation}
is closed isometric $(\|U_H x\| = \|x\|)$, and $(I - U_H)^{-1}$ exists. 
\end{theorem}


We have, moreover, 
\begin{equation}\label{eq:VII.4.5}
H = i (I + U _{H}) (I - U _{H}) ^{-1}.
\end{equation}

Thus, in particular, $D(H) = R(I - U_H)$ is dense in $X$. 

\begin{definition}\label{def:VII.4.Definition.3} $U_{H}$ is called the Cayley transform of $H$. 
\end{definition}

\begin{proof}[Proof of Theorem~\ref{thm:VII.4.Theorem.1}.]
We have, for any $x \in D(H)$,
\[
((H \pm i I) x, (H \pm i I) x) = (H x, H x) \pm (H x, i x) \pm (i x, H x) + (x, x).
\]

The symmetry condition for $H$ implies $(Hx, ix) = -i(Hx, x) = -i(x, Hx) = -(ix, Hx)$ and so 
\begin{equation}\label{eq:VII.4.6}
\left\| (H \pm i I) x \right\| ^{2} = \left\| H x \right\| ^{2} + \left\| x \right\| ^{2}.
\end{equation}

Hence $(H + iI)x = 0$ implies $x = 0$ and so the inverse $(H + iI)^{-1}$ exists. Since $\|(H + iI)x\| \geqq \|x\|$, the inverse $(H + iI)^{-1}$ is continuous. By \eqref{eq:VII.4.6}, it is clear that $\|U_Hy\| = \|y\|$, i.e., $U_H$ is isometric. 

$U_{H}$ is closed. For, let $(H + iI)x_{n} = y_{n}\rightarrow y$ and $(H - iI)x_{n} = z_{n}\rightarrow z$ as $n\to \infty$. Then we have, by \eqref{eq:VII.4.6}, $\| y_n - y_m\|^2 = \| H(x_n - x_m)\|^2 +\| x_n - x_m\|^2$, and so $(x_{n} - x_{m})\to 0$, $H(x_{n} - x_{m})\to 0$ as $n,m\to \infty$. Since $H$ is closed, we must have $x = s\text{- }\lim\limits_{n\to \infty}x_n\in D(H)$ and $s\text{- }\lim\limits_{n\to \infty}Hx_n = Hx$. Thus $(H + iI)x_{n}\to y = (H + iI)x,(H - iI)x_{n}\to z = (H - iI)x$ and so $U_Hy = z$. This proves that $U_{H}$ is closed. 

From $y = (H + iI)x$ and $U_Hy = (H - iI)x$, we obtain $2^{-1}(I - U_H)y = ix$ and $2^{-1}(I + U_H)y = Hx$. Thus $(I - U_H)y = 0$ implies $x = 0$ and so $(I + U_H)y = 2Hx = 0$ which implies $y = 2^{-1}((I - U_H)y + (I + U_H)y) = 0$. Therefore the inverse $(I - U_H)^{-1}$ exists. By the same calculation as above, we obtain 
\[
\begin{gathered}
H x = 2 ^{-1} (I + U _{H}) y = i (I + U _{H}) (I - U _{H}) ^{-1} x, \text{ that is, } H = i (I + U _{H}) (I - U _{H}) ^{-1}.
\end{gathered}
\]
\end{proof}

\begin{theorem}[J. \textsc{von} \textsc{Neumann} {\cite{neumann228}}]\label{thm:VII.4.Theorem.2} Let $U$ be a closed isometric operator such that $R(I - U)^{a} = X$. Then there exists a uniquely determined closed symmetric operator $H$ whose Cayley transform is $U$. 
\end{theorem}

\begin{proof}
 We first show that the inverse $(I-U)^{-1}$ exists. Suppose that $(I-U)y=0$. For any $z=(I-U)w\in R(I-U)$, we have, by the isometric property of $U$, $(y,w)=(Uy,Uw)$ as in Section~\hyperref[sec:VII.1]{$1$}. Hence 
\[
(y, z) = (y, w) - (y, U w) = (U y, U w) - (y, U w) = (U y - y, U w) = 0. \]

Hence, by the condition $R(I - U)^a = X$, $y$ must be $= 0$. Thus $(I - U)^{-1}$ exists. Put $H = i(I + U)(I - U)^{-1}$. Then $D(H) = D((I - U)^{-1}) =$ 
$R(I - U)$ is dense in $X$. We first prove that $H$ is symmetric. Let $x, y \in D(H) = R(I - U)$ and put $x = (I - U)u, y = (I - U)w$. Then $(Uu, Uw) = (u, w)$ implies that 
\[
\begin{aligned}
&(H x, y)  = (i (I + U) u, (I - U) w) = i ((U u, w) - (u, U w)) = ((I - U) u, i (I + U) w) = (x, H y).
\end{aligned}
\]

The proof of $U_H = U$ is obtained as follows. For $x = (I - U)u$, we have $Hx = i(I + U)u$ and so $(H + iI)x = 2iu$, $(H - iI)x = 2iUu$. Thus $D(U_H) = \{2iu; u \in D(U)\} = D(U)$, and $U_H(2iu) = 2i \cdot Uu = U(2iu)$. Hence $U_H = U$. 


To complete the proof of Theorem~\ref{thm:VII.4.Theorem.2}, we show that $H$ is a closed operator. In fact, $H$ is the operator which maps $(I - U) u$ onto $i(I + U) u$. If $(I - U) u_n$ and $i(I + U) u_n$ both converge as $n \to \infty$, then $u_n$ and $U u_n$ both converge as $n \to \infty$. Hence by the closure property of $U$, we must have 
\[
u _{n} \rightarrow u, (I - U) u _{n} \rightarrow (I - U) u, i (I + U) u _{n} \rightarrow i (I + U) u.
\]

This proves that $H$ is a closed operator.
\end{proof}

For the structure of the adjoint operator of a symmetric operator, we have 

\begin{theorem}[J. von Neumann {\cite{neumann228}}]\label{thm:VII.4.Theorem.3} Let $H$ be a closed symmetric operator in a Hilbert space $X$. For the Cayley transform $U_{H} = (H - iI)(H + iI)^{-1}$ of $H$, we set 
\begin{equation}\label{eq:VII.4.7}
X _{H} ^{+} = D (U _{H}) ^{\perp}, X _{H} ^{-} = R (U _{H}) ^{\perp}.
\end{equation}
\end{theorem}

Then we have 
\begin{equation}\label{eq:VII.4.8}
X _{H} ^{+} = \{x \in X; H ^{*} x = i x \}, X _{H} ^{-} = \{x \in X; H ^{*} x = - i x \},
\end{equation}
and the element $x$ of $D(H^{*})$ is uniquely expressed as 
\begin{equation}x = x _{0} + x _{1} + x _{2}, \text{ where } x _{0} \in D (H), x _{1} \in X _{H} ^{+}, x _{2} \in X _{H} ^{-} \text{ so that } H ^{*} x = H x _{0} + i x _{1} + (- i x _{2}). \label{eq:VII.4.9}\end{equation}

\begin{proof}
 $x \in D(U_H)^\perp = D((H + iI)^{-1})^\perp$ implies $(x, (H + iI)y) = 0$ for all $y \in D(H)$. Hence $(x, Hy) = -(x, iy) = (ix, y)$ and so $x \in D(H^*)$, $H^*x = ix$. The last condition implies $(x, (H + iI)y) = 0$ for all $y \in D(H)$, i.e. $x \in D((H + iI)^{-1})^\perp = D(U_H)^\perp$. This proves the first half of \eqref{eq:VII.4.8}; the latter half may be proved similarly. 

Since $U_{H}$ is a closed isometric operator, we see that $D(U_{H})$ and $R(U_{H})$ are closed linear subspaces of $X$. Hence any element $x \in X$ is uniquely decomposed as the sum of an element of $D(U_{H})$ and an element of $D(U_{H})^{\perp}$. If we apply this orthogonal decomposition to the element $(H^{*} + iI)x$, we obtain 
\[
(H ^{*} + i I) x = (H + i I) x _{0} + x ^{\prime} \quad \text{ where } \quad x _{0} \in D (H), x ^{\prime} \in D (U _{H}) ^{\perp}.
\]

But we have $(H + iI)x_0 = (H^* +iI)x_0$ by $x_0\in D(H)$ and $H\subseteq H^{*}$. We have also $H^{*}x^{\prime} = ix^{\prime}$ by $x^{\prime}\in D(U_{H})^{\perp}$ and \eqref{eq:VII.4.8}. Thus 
\[
x ^{\prime} = (H ^{*} + i I) x _{1}, x _{1} = (2 i) ^{-1} x ^{\prime} \in D (U _{H}) ^{\perp},
\]
and so 
\begin{equation}(H ^{*} + i I) x = (H ^{*} + i I) x _{0} + (H ^{*} + i I) x _{1} \text{ where } x _{0} \in D (H), x _{1} \in D (U _{H}) ^{\perp} \notag\end{equation}

Therefore $(x - x_0 - x_1) \in R(U_H)^\perp$ by $H^*(x - x_0 - x_1) = -i(x - x_0 - x_1)$ and \eqref{eq:VII.4.8}. This proves \eqref{eq:VII.4.9}. The uniqueness of the representation \eqref{eq:VII.4.9} is proved as follows. Let $0 = x_0 + x_1 + x_2$ with $x_0 \in D(H)$, $x_1 \in D(U_H)^\perp$, $x_2 \in R(U_H)^\perp$. Then, by $H^* x_0 = Hx_0$, $H^* x_1 = ix_1$, $H^* x_2 = -ix_2$, 
\[
0 = \left(H ^{*} + i I\right) 0 = \left(H ^{*} + i I\right) \left(x _{0} + x _{1} + x _{2}\right) = (H + i I) x _{0} + 2 i x _{1}.
\]

But by the uniqueness of the orthogonal decomposition of $X$ as the sum of $D(U_{H})$ and $D(U_{H})^{\perp}$, we obtain $(H + iI)x_{0} = 0$, $2ix_{1} = 0$. Since the inverse $(H + iI)^{-1}$ exists, we must have $x_{0} = 0$ and so $x_{2} = 0 - x_{0} - x_{1} = 0 - 0 - 0 = 0$. 
\end{proof}

\begin{corollaryu} $A$ closed symmetric operator $H$ in a Hilbert space $X$ is self-adjoint iff its Cayley transform $U_{H}$ is unitary. 
\end{corollaryu}

\begin{proof}
 The condition $D(H) = D(H^{*})$ is equivalent to the condition $D(U_{H})^{\perp} = R(U_{H})^{\perp} = \{0\}$. The last condition in turn is equivalent to the condition that $U_{H}$ is unitary, i.e. the condition that $U_{H}$ maps $X$ onto $X$ one-one and isometrically. 
\end{proof}


\yossection{5}{The Closed Range Theorem}
\label{sec:VII.5}

The closed range theorem of S. \textsc{Banach} \cite{banach9} reads as follows. 

\begin{theoremu} Let $X$ and $Y$ be $B$-spaces, and $T$ a closed linear operator defined in $X$ into $Y$ such that $D(T)^{a} = X$. Then the following propositions are all equivalent: 
\begin{equation}\label{eq:VII.5.1}
R(T) \text{ is closed in } Y,
\end{equation}
\end{theoremu}

\begin{align}
&\quad R \left(T ^{\prime}\right) \text{ is closed in } X ^{\prime}, \label{eq:VII.5.2} \\
R (T)& = N \left(T ^{\prime}\right) ^{\perp} = \{y \in Y; \langle y, y ^{*} \rangle = 0 \text{ for all } y ^{*} \in N \left(T ^{\prime}\right) \}, \label{eq:VII.5.3} \\
R (T ^{\prime})& = N (T) ^{\perp} = \left\{x ^{*} \in X ^{\prime}; \langle x, x ^{*} \rangle = 0 \text{ for all } x \in N (T) \right\}. \label{eq:VII.5.4}
\end{align}

\begin{proof}
 The proof of this theorem requires five steps. 


The first step. The proof of the equivalence $(1) \leftrightarrow \eqref{eq:VII.5.2}$ is reduced to the equivalence $(1) \leftrightarrow \eqref{eq:VII.5.2}$ for the special case when $T$ is a continuous linear operator such that $D(T) = X$. 

The graph $G = G(T)$ of $T$ is a closed linear subspace of $X \times Y$, and so $G$ is a $B$-space by the norm $\|\{x, y\}\| = \|x\| + \|y\|$ of $X \times Y$. Consider a continuous linear operator $S$ on $G$ into $Y$: 
\[
S \{x, T x \} = T x.
\]

Then the dual operator $S'$ of $S$ is a continuous linear operator on $Y'$ into $G'$, and we have 
\[
\begin{aligned}
&\langle \{x, T x \}, S ^{\prime} y ^{*} \rangle  = \langle S \{x, T x \}, y ^{*} \rangle = \langle T x, y ^{*} \rangle = \langle \{x, T x \}, \{0, y ^{*} \} \rangle , x \in D (T), y ^{*} \in Y ^{\prime}.
\end{aligned}
\]

Thus the functional $S' \cdot y^* - \{0, y^*\} \in (X \times Y)' = X' \times Y'$ vanishes at every point of $G$. But, the condition $\langle \{x, Tx\}, \{x^*, y_1^*\} \rangle = 0, x \in D(T)$, is equivalent to the condition $\langle x, x^* \rangle = \langle -Tx, y_1^* \rangle, x \in D(T)$, that is, to the condition $-T' y_1^* = x^*$. Hence 
\[
S ^{\prime} \cdot y ^{*} = \{0, y ^{*} \} + \{- T ^{\prime} y _{1} ^{*}, y _{1} ^{*} \} = \{- T ^{\prime} y _{1} ^{*}, y ^{*} + y _{1} ^{*} \}, y ^{*} \in Y ^{\prime}.
\]

By the arbitrariness of $y^*$, we see that $R(S') = R(-T') \times Y' = R(T') \times Y'$. Therefore $R(S')$ is closed in $X' \times Y'$ iff $R(T')$ is closed in $X'$, and, since $R(S) = R(T)$, $R(S)$ is closed in $Y$ iff $R(T)$ is closed in $Y$. Hence we have only to prove the equivalence \eqref{eq:VII.5.1} $\leftrightarrow$ \eqref{eq:VII.5.2} in the special case of a bounded linear operator $S$, instead of the original $T$. 

The second step. Let $X$ and $Y$ be $B$-spaces, and $T$ a bounded linear operator on $X$ into $Y$. Then $(1) \rightarrow \eqref{eq:VII.5.2}$. 

We consider $T$ as a bounded linear operator $T_{1}$ on $X$ into the $B$-space $Y_{1} = R(T)^{a} = R(T)$. We have to prove that \eqref{eq:VII.5.2} is true. $T_{1}'y_{1}^{*}, y_{1}^{*} \in Y_{1}'$, is defined by 
\[
\langle T _{1} x, y _{1} ^{*} \rangle = \langle T x, y _{1} ^{*} \rangle = \langle x, T _{1} ^{\prime} y _{1} ^{*} \rangle , x \in D (T _{1}) = D (T) = X.
\]

By the Hahn-Banach theorem, the functional $y_1^*$ can be extended to a $y^* \in Y'$ in such a way that $\langle Tx, y_1^* \rangle = \langle Tx, y^* \rangle$, $x \in D(T) = X$. Hence $T_1' y_1^* = T' y_1^*$ and so $R(T_1') = R(T')$. Thus it suffices to assume that $R(T) = Y$. Then, by the open mapping theorem in Chapter~\ref{chap:II}, \ref{sec:II.5}, there exists a $c > 0$ such that for each $y \in Y$, there exists an $x \in X$ with $Tx = y$, $\|x\| \leqq c \|y\|$. Thus, for each $y^*$ in $D(T')$, we have 
\[
| \langle y, y ^{*} \rangle | = | \langle T x, y ^{*} \rangle | = | \langle x, T ^{\prime} y ^{*} \rangle |
\]

Hence 
\begin{align}
&\leqq \left\| x \right\| \cdot \left\| T ^{\prime} y ^{*} \right\| \leqq c \left\| y \right\| \cdot \left\| T ^{\prime} y ^{*} \right\|. \notag \\
\left\| y ^{*} \right\|& = \sup\limits _{\left\| y \right\| \leqq 1} | \langle y, y ^{*} \rangle | \leqq c \| T ^{\prime} y ^{*} \| \notag
\end{align}
and so $(T')^{-1}$ exists and is continuous. Moreover, $(T')^{-1}$ is a closed linear operator as the inverse of a bounded linear operator. Hence we see that the domain $D((T')^{-1}) = R(T')$ must be closed in $X'$. 

The third step. Let $X$ and $Y$ be $B$-spaces, and $T$ a bounded linear operator on $X$ into $Y$. Then \eqref{eq:VII.5.2} $\to$ \eqref{eq:VII.5.1}. 

As in the second step, we consider $T$ as a bounded linear operator $T_{1}$ on $X$ into $Y_{1} = R(T)^{a}$. Then $T_{1}'$ has the inverse, since $T_{1}'y_{1}^{*} = 0$ implies 
\[
\langle T _{1} x, y _{1} ^{*} \rangle = \langle T x, y _{1} ^{*} \rangle = \langle x, T _{1} ^{\prime} y _{1} ^{*} \rangle = 0, x \in D (T _{1}) = D (T) = X,
\]
and so, since $R(T_1) = R(T)$ is dense in $Y_1 = R(T)^a$, $y_1^*$ must be $0$. Therefore, the condition that $R(T') = R(T_1')$ (proved above) is closed, implies that $T_1'$ is a continuous linear operator on the $B$-space $(R(T)^a)' = Y_1$ onto the $B$-space $R(T_1')$ in a one-one way. Hence, by the open mapping theorem, $(T_1')^{-1}$ is continuous. 

We then prove that $R(T)$ is closed. To this purpose, it suffices to derive a contradiction from the condition 
there exists a positive constant $\varepsilon$ such that the image $\{T _{1} x; \| x \| \leqq \varepsilon \}$ is not dense in all the spheres $\| y \| \leqq n ^{-1}   (n = 1, 2, \dots)$ of $Y _{1} = R (T) ^{a} = R (T _{1}) ^{a}.$


For, if otherwise, the proof of the open mapping theorem shows that $R(T_{1}) = R(T) = Y_{1}$. Thus we assume that there exists a sequence $\{y_{n}\} \subseteq Y_{1}$ with 
\[
\slim _{n \to \infty} y _{n} = 0, y _{n} \bar {\in} \{T _{1} x; \| x \| \leqq \varepsilon \} ^{a} (n = 1, 2, \dots).
\]

Since $\{T_{1}x;\|x\|\leqq\varepsilon\}^{a}$ is a closed convex, balanced set of the $B$-space $Y_{1}$ there exists, by Mazur's theorem in Chapter~\ref{chap:IV}, \ref{sec:IV.6}, a continuous linear functional $f_{n}$ on the $B$-space $Y_{1}$ such that 
\[
f _{n} (y _{n}) > \sup\limits _{\| x \| \leqq \varepsilon} | f _{n} (T _{1} x) | \quad (n = 1, 2, \dots).
\]

Hence $\| T_1'f_n\| < \varepsilon^{-1} \| f_n\| \| y_n\|$, and so, by $\slim\limits_{n\to \infty}y_n = 0$, $T_1'$ does not have a continuous inverse. This is a contradiction, and so $R(T)$ must be closed. 

The fourth step. We prove \eqref{eq:VII.5.1} $\rightarrow$ \eqref{eq:VII.5.3}. First, it is clear, from 
\[
\langle T x, y ^{*} \rangle = \langle x, T ^{\prime} y ^{*} \rangle , x \in D (T), y ^{*} \in D \left(T ^{\prime}\right),
\]
that $R(T) \subseteq N(T')^\perp$. We show that \eqref{eq:VII.5.1} implies $N(T')^\perp \subseteq R(T)$. Assume that there exists a $y_0 \in N(T')^\perp$ with $y_0 \bar{\in} R(T)$. Then, by the Hahn-Banach theorem, there exists a $y_0^* \in Y'$ such that $\langle y_0, y_0^* \rangle \neq 0$ and $\langle Tx, y_0^* \rangle = 0$ for all $x \in D(T)$. The latter condition implies $\langle x, T'y_0^* \rangle = 0$, $x \in D(T)$, and hence $T'y_0^* = 0$, i.e., $y_0 \bar{\in} N(T')^\perp$. This is a contradiction and so we must have $N(T')^\perp \subseteq R(T)$. 

The implication $\eqref{eq:VII.5.3} \to (1)$ is clear, since $N(T')^\perp$ is closed by virtue of the continuity in $y$ of $\langle y, y^* \rangle$. 

The fifth step. We prove \eqref{eq:VII.5.2} $\rightarrow$ \eqref{eq:VII.5.4}. The inclusion $R(T') \subseteq N(T)^\perp$ is clear as in the case of \eqref{eq:VII.5.3}. We show that \eqref{eq:VII.5.2} implies that $N(T)^\perp \subseteq R(T')$. To this purpose, let $x^* \in N(T)^\perp$, and define, for $y = Tx$, the functional $f_1(y)$ of $y$ through $f_1(y) = \langle x, x^* \rangle$. It is a one-valued function of $y$, since $Tx = Tx'$ implies $(x - x') \in N(T)$ and so, by $x^* \in N(T)^\perp$, $\langle (x - x'), x^* \rangle = 0$. Thus $f_1(y)$ is a linear functional of $y$. \eqref{eq:VII.5.2} implies \eqref{eq:VII.5.1}, and so, by the open mapping theorem applied to the operator $S$ in the first step, we may choose the solution $x$ of the equation $y = Tx$ in such a way that \slim $y = 0$ implies \slim $x = 0$. Hence $f_{1}(y) = \langle x, x^{*} \rangle$ is a continuous linear functional on $Y_{1} = R(T)$. Let $f \in Y'$ be an extension of $f_{1}$. Then 
\[
f (T x) = f _{1} (T x) = \langle x, x ^{*} \rangle .
\]

This proves that $T'f = x^*$. Hence $N(T)^\perp \subseteq R(T')$. 

That \eqref{eq:VII.5.4} implies \eqref{eq:VII.5.2} is clear, since $\langle x, x^* \rangle$ is a continuous linear functional of $x$.
\end{proof}

\begin{corollary}\label{cor:VII.5.Corollary.1} Let $X$ and $Y$ be $B$-spaces, and $T$ a closed linear operator on $D(T) \subseteq X$ into $Y$ such that $D(T)^a = X$. Then 
\begin{align}
R (T)& = Y \quad \text{ iff } \quad T ^{\prime} \text{ has a continuous inverse }, \label{eq:VII.5.5} \\
R (T ^{\prime})& = X ^{\prime} \text{ iff } T \text{ has a continuous inverse. } \label{eq:VII.5.6}
\end{align}
\end{corollary}

\begin{proof}
 Suppose that $R(T) = Y$. Then, from $\langle Tx, y^* \rangle = \langle x, T'y^* \rangle$, $x \in D(T)$ and $T'y^* = 0$, we obtain $y^* = 0$; that is, $T'$ must have the inverse $(T')^{-1}$. Since, by $R(T) = Y$ and \eqref{eq:VII.5.2}, $R(T')$ is closed, the closed graph theorem implies that $(T')^{-1}$ is continuous. Next let $T'$ admit a continuous inverse. Then $N(T') = \{0\}$ and also \eqref{eq:VII.5.2} holds since $T'$ is closed. Thus, by \eqref{eq:VII.5.3}, $R(T) = Y$. 

Suppose that $R(T') = X'$. Then, from $\langle Tx, y^* \rangle = \langle x, T'y^* \rangle$, $y^* \in D(T')$ and $Tx = 0$, we obtain $x = 0$, i.e., $T$ must have the inverse $T^{-1}$. Since, by $R(T') = X'$ and \eqref{eq:VII.5.1}, $R(T)$ is closed, the closed graph theorem implies that $T^{-1}$ must be continuous. Next let $T$ admit a continuous inverse. Then $N(T) = \{0\}$ and also \eqref{eq:VII.5.1} holds since $T$ is closed. Thus, by \eqref{eq:VII.5.4}, $R(T') = X'$. 
\end{proof}

\begin{corollary}\label{cor:VII.5.Corollary.2} Let $X$ be a Hilbert space with a scalar product $(u, v)$, and $T$ a closed linear operator with dense domain $D(T) \subseteq X$ and range $R(T) \subseteq X$. Suppose that there exists a positive constant $c$ such that 
\begin{equation}\label{eq:VII.5.7}
\operatorname{Re} (T u, u) \geqq c \| u \| ^{2} \quad \text{ for all } \quad u \in D (T).
\end{equation}
\end{corollary}

Then $R(T^{*}) = X$. 

\begin{proof}
 By Schwarz' inequality, we have 
\[
\left\| T u \right\| \cdot \left\| u \right\| \geqq \operatorname{Re} (T u, u) \geqq c \| u \| ^{2} \text{ for all } u \in D (T). \]

Hence $\|Tu\| \geqq c \|u\|, u \in D(T)$, and so $T$ admits a continuous inverse. Thus, by the preceding Corollary, $R(T') = X$. Hence $R(T^*) = R(T') = X$. 
\end{proof}

\begin{remarku} $A$ linear operator $T$ on $D(T) \subseteq X$ into $X$ is called accretive (the terminology is due to K. \textsc{Friedrichs} and T. \textsc{Kato}) if 
\begin{equation}\label{eq:VII.5.8}
\operatorname{Re} (T u, u) \geqq 0 \quad \text{ for all } \quad u \in D (T).
\end{equation}
\end{remarku}

$T$ is called dissipative (the terminology is due to R. S. \textsc{Phillips}) if $-T$ is accretive. 

\yossub{References for Chapter VII}

For a general account concerning Hilbert spaces, see M. H. \textsc{Stone} \cite{stone311}, N. I. \textsc{Achieser}-I. M. \textsc{Glazman} \cite{achieser1} and N. \textsc{Dunford}-J. \textsc{Schwartz} \cite{dunford48}. The closed range theorem is proved essentially in S. \textsc{Banach} \cite{banach9}. 
